\section{Ring-class interpolation and integral deformations}
\label{rl:section}

The ring-class argument has two outputs: a lower bound for the index of
a Heegner point at a known simple zero, and a transfer of nonvanishing to
an unknown base.  The first is uniform as the companion field varies
with bounded ramification; the second fixes that field before a binary
family grows.  We state both outputs before constructing their cochain
and evaluation maps.  Their determinant proofs occupy
Section~\ref{rd:section}.

Fix $E/\mathbf Q$ of conductor $N$, and write $T=T_2E$ and
$\Lambda=\mathbf Z_2[[t]]$.  The imaginary quadratic field
$K=\mathbf Q(\sqrt k)$ satisfies the split Heegner hypothesis at $2N$.
We omit the two fields with extra units and, when $E$ has complex
multiplication, its CM field.  Let $h$ be a product of distinct signed
good odd primes, all split in $K$, with $(h,k)=1$; $h=1$ is allowed.
The primes dividing $h$ are called the \emph{active conductor primes}.
The character $\chi=\chi_h$ is the corresponding quadratic character
over $\mathbf Q$.

Fix a modular parametrization $\pi:X_0(N)\to E$ with $\pi(\infty)=0$.
Write $H_c$ for the ring-class field of $K$ of conductor $c$, and $y_c$
for the image under $\pi$ of a conductor-$c$ CM point, with fixed CM
orientations.  The restriction of $\chi_h$ to $G_K$ factors through
$\operatorname{Gal}(H_{|h|}/K)$.  The primitive weighted Heegner sum is
\[
 P_h=\sum_{\sigma\in\operatorname{Gal}(H_{|h|}/K)}
             \chi_h(\sigma)\,\sigma(y_{|h|}).
\]
Through the quadratic-twist identification, it is a point of
$E^{(h)}(K)$.  When the product of $L(E^{(h)},s)$ and $L(E^{(hk)},s)$
has a simple zero,
the Heegner-point theorem makes $P_h$ nontorsion and
$E^{(h)}(K)$ has rank one.  Let $j_h$ be the $2$-adic index valuation of $P_h$
in the free rank-one lattice
$E^{(h)}(K)^\wedge_2/\mathrm{torsion}$.  The quantities $w(h)$,
$s(h)$, and $s(hk)$ have the meanings fixed in
Section~\ref{sec:preliminaries}; the two Shafarevich--Tate lengths are
finite in this simple-zero case.

\begin{corollary}[Ring-class lower bound]
\label{rd:ring-lower}
Every simple-zero ring instance in the stated data satisfies
\begin{equation}
 2j_h\ge 4w(h)+s(h)+s(hk)-C(E,A),
       \qquad \omega(k)\le A.
 \label{rd:ring-lower-equation}
\end{equation}
\end{corollary}

For the second output, fix $K$ and the base $h_0=1$.  A binary family
has the form
\begin{equation}
 h_x=\prod_{q\in Q}(q^*)^{\lambda_q(x)},\qquad
 \lambda_q:\mathbb F_2^b\longrightarrow\mathbb F_2\text{ linear},
 \quad x\in\mathbb F_2^b.
 \label{rd:binary-data}
\end{equation}
All primes in $Q$ are good, odd, disjoint from the fixed support, and
split in $K$.  All vertices have the same local squareclasses at $2N$
and the same required real sign.

\begin{theorem}[The fixed-\(K\) missing-vertex implication]
\label{rd:missing-vertex}
Fix the ring data \(E,K\), with
\begin{equation}
 \operatorname{corank}_{\mathbb Z_2}
       \operatorname{Sel}_{2^\infty}(E/K)=1.
 \label{rd:base-corank}
\end{equation}
Suppose binary families \eqref{rd:binary-data}
exist for arbitrarily large \(b\), and for
every nonzero \(x\) the product
\(L(E^{(h_x)},s)L(E^{(h_xk)},s)\) has a simple zero and
\begin{equation}
 2j_{h_x}-4w(h_x)-s(h_x)-s(h_xk)\le B
 \label{rd:nonzero-bound}
\end{equation}
with \(B\) independent of the family.  Then
\begin{equation}
 \operatorname{ord}_{s=1}
      \bigl(L(E,s)L(E^{(k)},s)\bigr)=1.
 \label{rd:missing-conclusion}
\end{equation}
This implication does not assume Mordell--Weil
rank one or finite Shafarevich--Tate group
at the base.
\end{theorem}

To prove these assertions, we compare the Heegner-point indices with
integral Selmer determinants.  The present section constructs the
limiting complexes and the maps needed to test their classes at new
primes.  Taking limits of differential matrices alone would not suffice:
a cycle in the limit may lift only to approximate cycles at finite
precision, and its localization must still be evaluation of a crossed
cocycle.

There are two passages to a limit.  In the inner construction $K$ and a
finite active support are fixed while an auxiliary scalar-conductor
prime $r_i$ and the coefficient precision vary.  Additional inert
\emph{derivative primes} will realize evaluations and switch local
conditions.  In the outer construction $K$ and the active support may
vary, but $\omega(k)$ and the number of retained derivative primes are
bounded.  This second passage gives the uniform lower bound; the final
fixed-$K$ transfer instead uses the bounded clearing factor of
Proposition~\ref{rd:family-clearing}.

We use
\[
 A_i=(\mathbf Z/2^i\mathbf Z)[t]/(t^i).
\]
The same constructions work over $A_i[G]$ and $\Lambda[G]$, where $G$ is a
fixed finite elementary abelian $2$-group and the additional character has
values in $G$.  Statements uniform in the number of generators of $G$ will
be identified separately; finite-diagram compactness by itself is a statement
for fixed $G$.

\subsection{Free models with maps to the original cochains}

\begin{lemma}[Finite Artin models]\label{rl:finite-model}
Let $(A,\mathfrak m,k)$ be an Artin local ring, and let $D$ be a complex of
flat $A$-modules.  Suppose $H^*(D\otimes_A k)$ is finite dimensional and
supported in $[a,b]$.  There is a bounded finite free complex $C$, with
\[
 \operatorname{rank}_A C^j=\dim_k H^j(D\otimes_A k),
\]
and a strong deformation retraction
\[
 i:C\longrightarrow D,\qquad p:D\longrightarrow C,\qquad
 pi=1,\qquad 1-ip=dh+hd.
\]
Thus maps to the original complex, as well as maps between finite models,
are available.  Any fixed finite collection of chain maps, homotopies,
triangles, and perfect pairings can be transported to these models.
\end{lemma}

\begin{proof}
Flat modules over an Artin local ring are free, including modules of infinite
rank.  Here is a useful version of the argument.  Lift a basis of $M/\mathfrak
mM$ to a flat module $M$, and map the corresponding free module to $M$.
Its cokernel is zero because $\mathfrak m$ is nilpotent.  Flatness shows that
the kernel has zero reduction, and nilpotence again makes the kernel zero.

Split the residual complex as its cohomology, with zero differential, plus
contractible two-term disks.  Lift all the graded bases.  These lifts are
graded isomorphisms: an inverse modulo $\mathfrak m$ lifts to an inverse by
a finite geometric series.  Let $d_0$ be the lifted disk differential and
write $d=d_0+\epsilon$.  All coefficients of $\epsilon$ belong to
$\mathfrak m$.  For the disk contraction $(i_0,p_0,h_0)$, arrange
$h_0^2=h_0i_0=p_0h_0=0$ and
$1-i_0p_0=d_0h_0+h_0d_0$.  The operators
$1+h_0\epsilon$ and $1+\epsilon h_0$ have finite geometric inverses.
The elementary perturbation identities give the differential
\[
 d_C=p_0\epsilon(1+h_0\epsilon)^{-1}i_0
\]
on the lifted residual cohomology, and give
\[
 i=(1+h_0\epsilon)^{-1}i_0,\quad
 p=p_0(1+\epsilon h_0)^{-1},\quad
 h=h_0(1+\epsilon h_0)^{-1}.
\]
Substitution, using $(d_0+\epsilon)^2=0$, proves the asserted contraction
identities.  These calculations are valid for infinite disk sums: each
operator is a module homomorphism and every geometric series terminates
at the nilpotence exponent of $\mathfrak m$.

For a map between original complexes use $p'fi$; insert the contraction
homotopies when composing maps.  A quasi-isomorphism between bounded free
complexes has a contractible cone and hence a homotopy inverse.  This gives
the stated finite-diagram assertion, including the homotopies expressing
triangles and perfect duality.
\end{proof}

For a profinite group and a finite free coefficient module over $A$,
continuous inhomogeneous cochains are flat: locally constant functions are
the filtered union of the modules of functions on finite clopen partitions,
and these modules are finite free.  Their reduction is the cochain complex
with reduced coefficients.  The same observations apply to restrictions,
mapping cones, and local conditions given by finite free complexes.  In the
applications below the groups are Galois groups of totally imaginary number
fields or nonarchimedean local fields.  Their $2$-cohomological dimension and
the usual finiteness theorems give the bounded residual cohomology required
in Lemma~\ref{rl:finite-model}.

\begin{lemma}[Limits of specified diagrams]\label{rl:finite-limits}
Suppose diagrams over $A_i$, or over $A_i[G]$, have uniformly bounded cohomological amplitudes and graded free
ranks, and include the chain homotopies and inverses for every asserted
identity or equivalence.  After taking a cofinal subsequence, or a
nonprincipal ultrafilter limit, their matrices define the same diagram of
perfect complexes over $\Lambda$, or $\Lambda[G]$.

Moreover, the contractions into the original cochains may be retained for
the evaluation construction in Lemma~\ref{rl:evaluation}.  Adding finitely
many new objects or maps does not require changing already chosen models
for the old objects.
\end{lemma}

\begin{proof}
There are finitely many possible graded ranks.  Fix them on an
ultrafilter-large set.  At a fixed precision $j$, all matrix entries belong
to the finite ring $A_j$ (or $A_j[G]$), so have unique ultrafilter limits.
These limits are compatible as $j$ varies.  Matrix identities, including
the homotopy-inverse identities, are finite polynomial identities and pass
to the limit.  New models and maps may be mapped to old models by their
already fixed contractions.  Passing to any slower cofinal precision has
the same old limits.  For the all-sequence evaluation statements below,
fix this ultrafilter once and for all, and keep the original contraction
at every stage.  Later arithmetic diagrams must extend these same
stages.  Restriction to an ultrafilter-large set and passage to slower
cofinal coefficient precision preserve every old evaluation on every
old group sequence.  An arbitrary later subsequence or a different
ultrafilter is not allowed in this assertion.  Notice that this argument preserves specified
arithmetic maps; it does not create an arithmetic map from a formal
description of its expected properties.
\end{proof}

\paragraph{Target-side data.}
The ring-class fields used to construct a class need not have bounded
degree.  Their cochain complexes are not objects of the limiting diagram.
At each finite stage, first perform descent, Shapiro transfer, the
coefficient weighting, and every coset sum.  These operations give
cochains on the fixed target field \(K\), with their local lifts.  Only
then apply Lemma~\ref{rl:finite-model} to the target global and local
complexes and their Selmer cones.

Here is the required inner rank bound.  Fix \(K\), the active support,
a finite set of derivative slots, and, when used, \(G\).  The permitted
ramification set \(S_i\) consists of this fixed support and a bounded
number of moving primes.  Modulo the maximal ideal of \(A_i\) or
\(A_i[G]\), all scalar characters are trivial and the representation is
\(E[2]\).  Put \(F=K(E[2])\), a fixed field.  Kummer theory bounds
\(H^1(G_{F,S_i},\mathbf F_2)\) by the rank of the \(S_i\)-unit group
and the dimension of the fixed group \(\operatorname{Cl}(F)[2]\).
Both bounds are uniform: the number of primes above \(S_i\) is bounded
and passing to the \(S_i\)-class group only decreases the required
class-group bound.  The fixed finite group
\(\operatorname{Gal}(F/K)\) and inflation--restriction then bound
\(H^1(G_{K,S_i},E[2])\).  Degree zero is bounded by two, and the global
Euler characteristic bounds degree two; higher degrees vanish because
\(K\) is totally imaginary.  Local cohomology has uniform bounds as
well: the local degrees over \(\mathbf Q\) are at most two, there are
only finitely many fixed local conditions at \(2N\), and the moving
good-prime and derivative complexes have fixed ranks.  The Selmer-cone
exact sequence gives the same conclusion for every target local
condition used below.  Thus Lemma~\ref{rl:finite-model} gives uniformly
bounded target ranks.  This inner bound may depend on the fixed active
support and is not a bound for a sequence of growing binary families.
The separate uniform bound after completion at \((2)\) is
Lemma~\ref{rl:residual-bound}; binary-family uniformity is obtained only
after the elimination in Proposition~\ref{rd:family-clearing}.

For clarity, let \(D_i\) be a target cochain complex and write its
contraction as \((\iota_i,p_i,h_i)\).  A transferred cycle \(z_i\)
is retained as the vector \(p_i z_i\).  Its representative in the
original target cochains differs from \(z_i\) by the boundary
\(d h_i z_i\).  If an arithmetic identity is
\(z_i-z'_i=d b_i\), retain the target vectors
\(p_i z_i,p_i z'_i,p_i b_i\); their identity is exact.  For a target
comparison map \(f_i:D_i\to D'_i\), retain
\(p'_i f_i\iota_i\), together with the contraction homotopy.  For
example,
\[
 p'_i f_i z_i-p'_i f_i\iota_i p_i z_i
       =d\bigl(p'_i f_i h_i z_i\bigr).
\]
Apply this also to local lifts in the Selmer cone, specialization
comparisons, and the homotopies expressing duality and local switches.
Each specified cycle or boundary witness is a vector in a bounded-rank
target term, so adds only finitely many coordinates.  Boundary witnesses
are compressed after transfer and have target rank; no list indexed by
the ring-class cosets is retained.  The lemma therefore retains all
these identities without retaining a source cochain complex or a
growing transfer matrix.  The original target contractions remain
available in degrees zero, one, and two for all-sequence evaluation.
All operations use the same original stages and ultrafilter; a loss of
precision is accommodated by one common slower cofinal precision for
the entire finite diagram.  Adding derivative slots means extending
each old stage by finitely many target objects, not taking a limit with
an unbounded number of objects.

\subsection{Kummer conditions, including their derived degree zero}

\begin{lemma}[Integral Kummer local condition]\label{rl:kummer}
Let $F$ be a finite extension of a nonarchimedean local field of
characteristic zero, and let $B/F$ be any elliptic curve.  Put
$M_F=B(F)^\wedge_2=\varprojlim_n B(F)/2^nB(F)$.  There is a perfect local
condition
\[
 U_F=M_F[-1]\longrightarrow R\Gamma(F,T_2B).
\]
For every $n\geq1$, derived reduction has
\[
 H^0(U_F\otimes^{\mathbf L}\mathbf Z/2^n)=B(F)[2^n],\qquad
 H^1(U_F\otimes^{\mathbf L}\mathbf Z/2^n)=B(F)/2^nB(F).
\]
Its map in degree zero is the isomorphism onto
$H^0(F,B[2^n])$, and its map in degree one is the finite Kummer injection.
The condition is its own exact orthogonal complement under local Tate
duality.  These assertions hold at $F\mid2$, for all reduction types, and
without a condition on $B[2]$.
\end{lemma}

\begin{proof}
The local analytic description of $B(F)$ shows that $M_F$ is a finitely
generated $\mathbf Z_2$-module.  Its torsion is the finite group
$B(F)[2^\infty]$.  Consequently $M_F$ has a free resolution of length at
most one.  Also $H^0(F,T_2B)=0$, since $B(F)[2^\infty]$ has bounded
exponent.  The integral Kummer injection into $H^1(F,T_2B)$ therefore
defines the displayed morphism using the truncation in degrees at most
one of $R\Gamma(F,T_2B)$.

Tensor the free resolution of $M_F$ with $\mathbf Z/2^n$.  Its two
cohomology groups are $M_F[2^n]$ and $M_F/2^nM_F$.  The inverse limits of
the ordinary Kummer exact sequences identify their maps with the two
maps in the statement.  In particular finite torsion invariants have not
been discarded.

For orthogonality, compose cup product and the Weil pairing with the local
invariant map.  A morphism
\[
 M_F[-1]\otimes^{\mathbf L}M_F[-1]
       \longrightarrow\mathbf Z_2[-2]
\]
is determined by its homomorphism $M_F\otimes M_F\to\mathbf Z_2$:
the possible $\operatorname{Tor}_1$ is in degree one and gives no extra
degree-zero morphism to this target.  The indicated homomorphism vanishes,
because the finite Kummer image annihilates itself at every level by local
Tate duality \cite[I, Theorem~3.2 and Corollary~3.4]{MilneADT2006}.
Choose a nullhomotopy in free models.  The induced morphism
\[
 R\Gamma(F,T_2B)/U_F
       \longrightarrow R\operatorname{Hom}_{\mathbf Z_2}
                         (U_F,\mathbf Z_2)[-2]
\]
is a quasi-isomorphism.  This can be checked after reduction modulo $2$:
degree one is the perfect pairing of the finite Kummer quotient with the
Kummer image, and degree two is dual to $B(F)[2]$.  Its perfect cone is
therefore zero by derived Nakayama.  This proves exact orthogonality and
also its compatibility with derived reduction.  The derived duality of
the ambient local cochains is \cite[Proposition~5.2.4]{Nekovar2006};
the Kummer condition and its orthogonality have been constructed above.
No part of the proof uses an ordinary filtration or a classification of
reduction types.
\end{proof}

Where the scalar deformation is trivial, extend this local condition by
derived scalar extension to $A_i$ or $A_i[G]$.  For a finite free
coefficient representation $M$ unramified at a good prime $v$, use the
inflated residue-field complex
\[
 U_v=[\,M\xrightarrow{\operatorname{Fr}_v-1}M\,]
 \quad\text{in degrees }0,1.
\]
Its quotient in full local cochains is the singular local complex.
The full local condition is the identity on local cochains, and its
orthogonal strict condition is zero.  The conditions at a derivative
prime are different: once its local cochains have the split model
constructed in Lemma~\ref{rd:local-model}, they are the truncations
specified below.  Write that model as
\[
 L_v=D^0\oplus(F_v\oplus S_v)[-1]\oplus D^2[-2],
\]
where the differential is zero, $D^0$ is dual to $D^2$, and the finite
and transverse planes $F_v,S_v$ are each isotropic.  Put
\[
 U_{\rm str}=D^0,\qquad
 U_{\rm rel}=D^0\oplus(F_v\oplus S_v)[-1],\qquad
 U_F=D^0\oplus F_v[-1],\qquad U_S=D^0\oplus S_v[-1].
\]
Then $U_{\rm str}^{\perp}=U_{\rm rel}$ and each plane condition is
its own orthogonal complement.  Every displayed condition maps
isomorphically onto local $H^0$ and injectively into local $H^1$.
Thus they satisfy precisely the detection hypotheses used below.
Lemma~\ref{rd:local-model} proves their finite arithmetic realization
and supplies the isotropic-plane nullhomotopies.  These truncated
conditions are not identified with the full local condition.

For a finite permitted set $S$ the Selmer complex is, by definition,
\begin{equation}\label{rl:selmer-definition}
 C_S=\operatorname{Cone}\left(
 R\Gamma(G_{K,S},M)\oplus\bigoplus_{v\in S}U_v
 \longrightarrow\bigoplus_{v\in S}R\Gamma(K_v,M)
 \right)[-1].
\end{equation}
At $2N$ take the Kummer conditions just constructed.  Restriction and
localization triangles are the triangles of this cone.  The finite rings
$A_i$ and $A_i[G]$ are Frobenius coefficient rings; alternatively the
duality maps may be checked after reduction to their residue field.
Exact local orthogonality and global duality therefore give
\begin{equation}\label{rl:global-duality}
 C_S\simeq R\operatorname{Hom}(C_S^{\perp},A_i)[-3],
\end{equation}
with the appropriate dual representation.  This application uses the
orthogonality hypotheses of
\cite[Theorem~6.3.4]{Nekovar2006}; the theorem is not being used to assert
the existence of a varying-prime limit.  Complex places contribute zero
Tate complexes.  If the character is inverted by the nontrivial
automorphism of $K/\mathbf Q$, conjugate transport followed by the Weil
pairing identifies the two representations over the same scalar ring.
All maps and nullhomotopies in this paragraph are finite diagrams to
which Lemma~\ref{rl:finite-limits} applies.

\subsection{The moving conductor character}

\begin{proposition}[Finite-precision ring-class characters]\label{rl:characters}
Fix finitely many split primes of $K$ above odd rational primes, outside
$2N$, and orient one prime in each conjugate pair.  There are distinct
split rational primes $r_i$ avoiding all the fixed data and characters
\[
 \xi_i:\operatorname{Gal}(H_{r_i}/K)\longrightarrow
                 \mathbf Z/2^{n_i}\mathbf Z,
 \qquad n_i\longrightarrow\infty,
\]
such that $\xi_i$ is surjective on the inertia quotient at $r_i$, is zero
at the primes above $2N$, and has odd value at each prescribed oriented
prime.  We can require $r_i\to1$ in $\mathbf Z_2$ and
\begin{equation}\label{rl:trace-bound}
 v_2(2+a_{r_i}(E)),\quad v_2(2-a_{r_i}(E)),\quad
 v_2\det(1\pm\operatorname{Fr}_{r_i}\mid T)\ \leq c_E.
\end{equation}
The constant $c_E$ is independent of the list and of $K$.
\end{proposition}

\begin{proof}
We give the compatibility argument before imposing the Tate condition.
Choose prime ideals $\mathfrak p_1,\ldots,\mathfrak p_s$ with pairwise
disjoint conjugate pairs, including the oriented prescribed primes and
the primes above $2N$, and enough additional split prime ideals to
generate $\operatorname{Cl}(K)$.  Such additional ideals may avoid any
fixed finite set.  Put weight $w_j=1$ at a prescribed odd prime and
$w_j=0$ at every additional prime and every prime above $2N$.
The conjugate weight is $-w_j$.  Let
\[
 R=\ker\bigl(\mathbf Z^s\longrightarrow\operatorname{Cl}(K)\bigr),
\]
choose a basis $m^{(a)}$ of $R$, and choose $\beta_a\in K^\times$ with
$(\beta_a)=\prod_j\mathfrak p_j^{m_j^{(a)}}$.

Work first over $F_\infty=K(\mu_{2^\infty})$ and adjoin $2^n$-th roots
of the finitely many $\beta_a/\overline\beta_a$.  Prescribe on these
roots the multipliers
\begin{equation}\label{rl:kummer-prescription}
 \zeta_{2^n}^{\,-\sum_j m_j^{(a)}w_j}.
\end{equation}
The minus sign fixes the choice of connecting map in the ring-class
sequence and could be reversed throughout.  These prescriptions respect
every Kummer relation.  Indeed, if
\[
 \prod_a(\beta_a/\overline\beta_a)^{b_a}
                         \in(F_\infty^\times)^{2^n},
\]
take its valuation at a prime above each $\mathfrak p_j$ of nonzero
weight.  These odd primes are unramified in $F_\infty/K$.  Thus
$2^n$ divides $\sum_a b_am_j^{(a)}$ for each such $j$, proving the
required compatibility of~\eqref{rl:kummer-prescription}.  Primes over
$2$ have weight zero and impose no division by a ramification index.
Units impose no additional weighted relation.  Kummer duality supplies
the prescribed automorphism of the radical extension over $F_\infty$.

At a rational prime $r$ split in a sufficiently large finite cyclotomic
subfield, the power-residue character of $\mathbf F_r^\times$ is
surjective onto $\mu_{2^n}$.  The split ring-class exact sequence is
\[
 1\longrightarrow
  (\mathcal O_K/r\mathcal O_K)^\times/(\mathbf Z/r\mathbf Z)^\times
 \longrightarrow\operatorname{Pic}(\mathbf Z+r\mathcal O_K)
 \longrightarrow\operatorname{Cl}(K)\longrightarrow1.
\]
There is no unit correction under our exclusion of the extra-unit
fields.  The first group is $\mathbf F_r^\times$, with a principal
relation measured by $\beta/\overline\beta$ at the oriented prime over
$r$.  Consequently~\eqref{rl:kummer-prescription} says exactly that its
primitive power-residue character extends to the class group with the
specified values $w_j$.  Since the ring-class extension is unramified at
$2N$, value zero on Frobenius means that the resulting local character
is trivial there.  Conjugation acts by inversion.

We now show that imposing these finitely many radical conditions cannot
force either forbidden Tate trace.  In the non-CM case the image of
$G_{K(\mu_{2^\infty})}$ is an open subgroup of
$\operatorname{SL}_2(\mathbf Z_2)$, by the open-image theorem
\cite{Serre1972}.  Its index in the corresponding fixed-$E$ image changes
by at most two when $K$ changes.  Its commutator subgroup contains a
fixed open subgroup $U_E$ of $\operatorname{SL}_2(\mathbf Z_2)$.
One may see uniformity by starting with a principal congruence subgroup
in the fixed image: every index-two subgroup contains its squares, and
the commutators of the resulting fixed congruence subgroup are open.
Because the radical extension over $F_\infty$ is abelian, these
commutators act trivially on it.  Hence every permitted Tate coset
contains a translate of $U_E$.

For completeness, in the CM case let $M$ be the CM field.  The image over
$KM(\mu_{2^\infty})$ contains a fixed open subgroup of the norm-one
Cartan; the index changes by a bounded factor, as $[KM:\mathbf Q]\leq4$.
Choose $\gamma\in G_{KM}$ with
$\chi_{\rm cyc}(\gamma)=5^4=625$.  This is possible uniformly: the
cyclotomic image of $G_{KM}$ has index at most four in
$\mathbf Z_2^\times$, so contains fourth powers.  On the radical Galois
group, conjugation by $\gamma$ is raising to $625$, because all the
radicands lie in $K$.  On the Tate image over $KM$, conjugation is
trivial.  Therefore
\[
 \gamma u\gamma^{-1}u^{-625},\qquad
 u\in G_{KM(\mu_{2^\infty})},
\]
fixes every prescribed radical and acts on $T$ as $\rho_T(u)^{-624}$.
Raising to the fixed nonzero power $624$ carries an open subgroup of the
one-dimensional norm-one Cartan onto an open subgroup.  This proves the
required uniform kernel statement also in the CM case.  The open Cartan
image used here is the CM part of the usual $\ell$-adic image theorem
\cite[Chapter~IV, \S2.2, Remark]{Serre1968}.

On a coset of an open subgroup of $\operatorname{SL}_2$, trace is not
identically $2$ or $-2$.  The same holds on a coset of an open norm-one
Cartan subgroup; on the other Cartan-normalizer component trace is zero.
The relevant determinant-one sets are compact and have finitely many
cosets modulo the fixed open subgroup.  Choose in each coset an element
whose trace is neither forbidden value.  The largest of the finitely
many resulting valuations is a constant depending only on $E$.

Finally restrict the prescribed radical automorphism and such a Tate
lift to a finite Galois extension, including as many roots of unity as
required.  Chebotarev gives infinitely many rational primes with this
Frobenius class; choose the place over the prime to realize the chosen
representative.  They split in $K$ and have arbitrarily high prescribed
congruence to $1$.  Since
$\det(1\pm\operatorname{Fr}_r)=1\pm a_r+r$, a congruence precision
larger than the fixed trace bound gives all of~\eqref{rl:trace-bound}.
This is a finite Chebotarev condition at each stage, not a demand that
one prime split in an infinite extension.
\end{proof}

Choose $n_i$ large enough that $1+t$ has order dividing $2^{n_i}$ in
$A_i^\times$, and put $\Psi_i=(1+t)^{\xi_i}$.  For example $n_i\geq2i$
is sufficient.  After the matrix limits, the inertia exponent at the
moving prime is a unit of $\mathbf Z_2$, and a prescribed oriented
Frobenius exponent is odd.  We can orient the generator of inertia so
its action is exactly $1+t$.

\begin{lemma}[Local height-one tests]\label{rl:local-acyclic}
For the individual coefficient system $T\chi_h\Psi$, use full local
conditions at the moving prime $r$ and the active primes dividing $h$,
without adding primes from a larger binary family.  The limiting full
local complexes are acyclic over every height-one residue field of
$\Lambda$ of characteristic zero.  At $(2)$, with residue field
$\Omega=\mathbf F_2((t))$, the $r$-complex is acyclic, and so are the
complexes at active split primes whose prescribed exponent is odd.
These full conditions therefore agree with their orthogonal conditions
on the stated tests.
\end{lemma}

\begin{proof}
Use tame inertia and unramified Frobenius, or the Hochschild--Serre
complex for these two procyclic directions.  At $r$, inertia minus one
is multiplication by $t$ times a unit.  It is invertible on all the
indicated fields except the characteristic-zero fiber $t=0$.  There,
$r\to1$, and the two remaining Frobenius differentials have nonzero
determinant by~\eqref{rl:trace-bound}, including either sign of the
quadratic twist.  The bound also ensures that the integral local
torsion at this central fiber has bounded length.

At an active odd conductor prime, inertia of the quadratic character
acts by $-1$.  Its differential is thus $-2$ and is invertible in
characteristic zero.  In characteristic two the quadratic character
disappears.  At a split prime the Frobenius action is $zB$, where
$B\in\operatorname{GL}_2(\mathbf F_2)$ and
$z=(1+t)^a$ has odd $a\in\mathbf Z_2$.  Its coefficient of $t$ is one,
so $z$ is nonconstant.  No nonzero polynomial over the finite constant
field vanishes at $z$.  Thus both Frobenius complexes, with operators
$zB-1$ and $q^{-1}zB-1$, are invertible over $\Omega$; here $q=1$ in
the residue field.  This argument applies also to a reducible residual
representation.

The orthogonal full and strict conditions agree when full local cochains
are acyclic.  Equation~\eqref{rl:global-duality} then gives self-duality.
The global degree-zero assertion will follow from the evaluation maps
in Lemma~\ref{rl:evaluation}.
\end{proof}

\begin{lemma}[Uniform residual sizes]\label{rl:residual-bound}
Suppose $\omega(\operatorname{disc}K)\leq A$, the local twisting classes
at $2N$ belong to a fixed finite list, and all active conductor primes
are split with odd exponent.  After localization and completion at
$(2)$, the minimal free ranks of the Selmer complexes are bounded in
terms of $E,A$, and the number of retained derivative primes.  They do
not depend on the number of active conductor primes.
\end{lemma}

\begin{proof}
Reduce first over $\mathbf F_2[[t]]$.  The quadratic character is now
trivial.  Removing the active conductor primes from the permitted
ramification set changes the global complex by their singular local
complexes.  These, and their unramified local complexes, are acyclic
after inverting $t$, by Lemma~\ref{rl:local-acyclic}.  The resulting
smaller-support complex, including its comparison maps, can be
constructed at the original Artin stages before taking the limit.
Its support consists of $2N,r$, and the retained derivative primes.
Its $\mathbf F_2((t))$ cohomology dimensions are bounded by those at
$t=0$: choose a minimal finite free model over $\mathbf F_2[[t]]$.

Here is a uniform bound for the latter residual cohomology.  Put
$F=\mathbf Q(E[2])$, a fixed field.  Then $FK/F$ is at most quadratic.
Enlarge the allowed set of primes of $F$ by those ramified in $K$.
Its cardinality is bounded by a constant depending on $E,A$ and the
number of derivative primes.  Kummer theory bounds the number of
generators of the maximal pro-$2$ Galois group over $F$ unramified
outside this set, using the fixed unit group, the fixed $2$-class group,
and the allowed valuations.  The subgroup defining $FK$ has index at
most two; the elementary Schreier bound bounds its generator number.
Every finite Galois $2$-extension of $FK$ with this ramification has
Galois closure of $2$-power degree over $F$, so this subgroup indeed
controls the required pro-$2$ extensions of $FK$.

Over $FK$ the residual Tate module is trivial.  Restriction and
inflation--restriction therefore bound the residual $H^1$ over $K$.
The $H^0$ bound is immediate, and the global Euler characteristic or
global duality bounds $H^2$.  The local terms at the remaining bounded
set have bounded dimensions, with degrees of the local fields at most
two over the fixed rational local fields.  At $2N$ there are only
finitely many local quadratic twisting classes.  The same bounds
apply to the Selmer cone.  Lemma~\ref{rl:finite-model} now bounds the
minimal ranks.  No bound on a torsion exponent is inferred from these
rank bounds.
\end{proof}

\subsection{Evaluation on every sequence, rather than only a chosen list}

\begin{lemma}[All-sequence evaluation]\label{rl:evaluation}
Fix $K$ and a sequence of the finite cochain constructions above.  Keep
the contraction maps into the original global cochains.  Let
$\Gamma_K=\prod_i G_K$ and
$\Gamma=\Gamma_K\rtimes\langle\tau\rangle\subset\prod_iG_{\mathbf Q}$,
with simultaneous complex conjugation $\tau$.  Thus $\Gamma_K$ has
index two in $\Gamma$; the larger unrestricted product of rational
Galois groups is not used for induction.  The matrix limit defines a
coefficient action of $\Gamma_K$ on $\Lambda^2$.  For every field $L$
receiving a map from $\Lambda$, there is a natural injection
\[
 H^1(C\otimes_\Lambda L)\longrightarrow
 H^1_{\rm abs}(\Gamma_K,L^2),
\]
where the right side denotes abstract crossed-homomorphism classes.
Degree-zero evaluation likewise injects
$H^0(C\otimes_\Lambda L)$ into $(L^2)^{\Gamma_K}$.
Here $C$ may be the global complex or a Selmer complex whose local
condition maps induce isomorphisms on local $H^0$ and injections on
local $H^1$.  The same statement holds for the induced representation
of $\Gamma$.

There are evaluation maps on arbitrary sequences, not just on the
finite list used to detect cohomology.  Their degree-two identities
remain valid after arbitrary coefficient extension.
\end{lemma}

\begin{proof}
Let $C_i$ be a finite model and let $\iota_i$ be its chosen map to the
original global cochains, using the projection from the Selmer cone
when appropriate.  For every $g=(g_i)$ and $h=(h_i)$ define matrices
\[
 e^0_i:C_i^0\to A_i^2,\quad
 e^1_{g,i}:C_i^1\to A_i^2,\quad
 e^2_{g,h,i}:C_i^2\to A_i^2
\]
by evaluating $\iota_i$ in degrees zero, one, and two.  Their source
ranks are fixed and their entries belong to $A_i$.  Thus each has a
coefficientwise ultrafilter limit, for every choice of $g,h$.
No diagonal selection over the set of group elements is necessary:
one fixed ultrafilter defines all these limits.  Likewise
$\rho(g)=\lim_i\rho_i(g_i)$ is a representation.  The original
cochain identities imply
\begin{align}
 e^1_g d&=(\rho(g)-1)e^0,\label{rl:eval-boundary}\\
 e^2_{g,h}d&=\rho(g)e^1_h-e^1_{gh}+e^1_g.
                         \label{rl:eval-cocycle}
\end{align}
All identities pass to the limit because their matrix expressions
have finitely many terms.  A cycle $x\in C^1\otimes L$ consequently
gives the crossed homomorphism $g\mapsto e^1_g(x)$, and a boundary
gives a coboundary.

We prove injectivity in a way that also controls base change.  At each
finite stage let $k=\mathbf F_2$ be the residual field.  Choose a
space of representatives for residual $H^1$.  Evaluate pairs $(c,m)$,
where $c$ is in this space and $m\in k^2$, by
$b_c(g)-(\rho(g)-1)m$.  The intersection of the kernels of all
these evaluations consists precisely of pairs $(0,m)$ with $m$ an
invariant.  Since the vector space is finite dimensional, a list of
at most $\dim H^1+2$ elements already has this intersection.  For
these elements evaluation is a chain map
\[
 C_i\longrightarrow D_i=
 [\,A_i^2\xrightarrow{((\rho_i(g_j)-1))_j}(A_i^2)^d\,]
\]
in degrees zero and one, inducing an isomorphism on residual $H^0$
and an injection on residual $H^1$.  For Selmer complexes this uses
the exact sequence of~\eqref{rl:selmer-definition}: the indicated
local degree-zero isomorphisms give $H^0(C_i)=H^0_{\rm glob}$ and
$H^1(C_i)\hookrightarrow H^1_{\rm glob}$ residually.

Retain these matrices and their chosen group-element sequences.  The
cone of the limiting map $C\to D$ has zero residual cohomology in
degrees at most zero.  Successively cancel its unit differential
blocks.  The resulting bounded free cone starts in degree one.
After every coefficient extension it still starts in degree one;
the long exact sequence therefore makes $H^1(C\otimes L)\to
H^1(D\otimes L)$ injective.  A crossed cocycle that is a coboundary
on all of $\Gamma_K$ is in particular a coboundary on the detecting
list, proving the desired injection.  The same cone gives an
isomorphism on $H^0$ after coefficient extension.  Since $H^0(D\otimes L)$
is a subspace of $L^2$, degree-zero evaluation is injective; by
\eqref{rl:eval-boundary}, its image is fixed by every element of
$\Gamma_K$, not merely by the detecting list.  For induction the same proof
has coefficient rank four and the detecting bound is $\dim H^1+4$.
\end{proof}

We can now exclude global invariants on the tests of
Lemma~\ref{rl:local-acyclic}.  In characteristic two the nontrivial
$r$-inertia scalar excludes coefficient invariants.  In characteristic
zero the same argument works unless that scalar is trivial.  In the
latter case $t=0$, and $T\otimes\mathbf Q_2$ over $G_K$ is absolutely
irreducible: for non-CM curves this follows from the open image, and
for CM curves from the open Cartan and its nontrivial normalizer
component, since $K$ is not the CM field.  Twisting by $\chi_h$ does
not alter its invariant subspaces.  The degree-zero injection just
proved therefore gives $H^0=0$ for the limiting Selmer complex.
Together with \eqref{rl:global-duality}, this leaves degrees one and
two paired and kills degree three as well.

\begin{lemma}[The exact common kernel]\label{rl:kernel}
In the notation of Lemma~\ref{rl:evaluation}, let $H$ be the group of
sequences each of whose entries acts trivially on the full Tate module,
on all ring-class fields of $K$, and on the quadratic characters in
use.  Restriction from the crossed classes in that Lemma to $H$ is
injective at every characteristic-zero test.  It is also injective
at the residual $(2)$ test.  The characteristic-zero annihilator
needed for this assertion can be chosen with a fixed finite
$2$-valuation depending only on $E$, also as $K$ varies.
\end{lemma}

\begin{proof}
The Tate image contains an open subgroup of scalar homotheties, both
in the CM and non-CM cases
\cite[\S2, Theorem~3 and paragraph~7]{Bogomolov1981Points}.  Choose once and for all
$g\in G_{\mathbf Q}$ with Tate action $u\,1$, where
$u\in1+4\mathbf Z_2$ and $u\ne1$.  Its square belongs to every
quadratic subgroup.  Put
\[
 z=g^2\tau g^2\tau^{-1}.
\]
The generalized dihedral property of every ring-class extension
implies that $z$ is ring-class trivial.  It is also trivial on every
quadratic character, and its Tate action is the fixed scalar $u^4$.
For every sequence in $\Gamma_K$, its commutator with the constant
sequence $z$ is in the exact group $H$, since its Tate commutator and
all its abelian character images are trivial.  Thus $z$ is central in
$\Gamma_K/H$ and acts on the coefficients by $u^4$.

Here the elementary central-scalar argument suffices.  If a crossed
cocycle $b$ is defined on a group with such a central element, comparison
of $b(zg)$ and $b(gz)$ gives
\[
 (u^4-1)b(g)=(g-1)b(z).
\]
In characteristic zero this makes $b$ a coboundary.  Inflation--restriction
then gives injectivity on $H$.  The valuation of $u^4-1$ is fixed.

In characteristic two use instead, at the $i$-th stage, an inertia
element at $r_i$ mapping to the chosen generator of $\xi_i$.  Its
Tate action is trivial and its ring-class image commutes with all
ring-class images.  The associated sequence is therefore central
modulo $H$, and its coefficient action is $1+t$.  Since $t$ is
invertible in $\mathbf F_2((t))$, the same calculation applies.
It is essential here that $H$ is defined by the exact actions at each
stage, not merely by an action converging to the identity.
\end{proof}

\subsection{The outer characteristic-zero test}

\begin{proposition}[Integral evaluation through the outer limit]
\label{rl:outer-evaluation}
Let $C_j$ be any family of the preceding limiting complexes, completed
over $\mathcal O_j=\widehat{\Lambda_{(2)}}$, possibly for different
quadratic fields $K_j$.  Suppose their minimal ranks are bounded.
Keep their actual evaluation matrices, and transport them through the
integral contractions to these minimal models.  For a nonprincipal
ultrafilter $\mathcal U$, put
\[
 \mathcal O_\infty=
 \left(\prod_{\mathcal U}\mathcal O_j\right)/\mathfrak a,
 \qquad
 \mathfrak a=\{(x_j):
    \{j:v_2(x_j)\geq n\}\in\mathcal U\text{ for every }n\}.
\]
Then $\mathcal O_\infty$ is a DVR with uniformizer $2$ and a
characteristic-zero fraction field $L_\infty$.  The outer complex has
an injection of its $H^1$ over $L_\infty$ into crossed classes on the
product of the inner sequence groups.  Its restriction to the product
of their exact common kernels is injective.  All arbitrary-sequence
degree-two evaluation identities survive.  These conclusions also
hold after retaining any fixed finite collection of additional
integral diagrams.
\end{proposition}

\begin{proof}
A nonzero element of the quotient has a representative whose
valuations are bounded on an ultrafilter-large set.  A finite
partition then makes its valuation equal to a single integer $n$
on such a set.  It is $2^n$ times a unit.  Products add these finite
valuations, so the quotient is a domain, and every nonzero ideal has
an element of smallest valuation.  It is consequently a DVR.  The
integer $2^n$ is nonzero for every $n$, so its fraction field has
characteristic zero.

For each $j$, all-sequence evaluation maps originally have entries
in $\Lambda$.  Extension to $\mathcal O_j$ and composition with an
integral contraction still give integral matrices.  Thus evaluation
on any sequence of inner group elements has an entrywise image in
$\mathcal O_\infty$.  Equations~\eqref{rl:eval-boundary} and
\eqref{rl:eval-cocycle} survive for all such sequences.  In particular,
if an outer cycle is represented by vectors $x_j$ for which
$d_jx_j$ has valuation tending to infinity, the evaluation of
$d_jx_j$ does too: the degree-two evaluation matrix is integral.
This proves the crossed identity for an outer cycle without replacing
the $x_j$ by genuine inner cycles.

Over the inner residue field $\Omega_j$, choose a detecting list as
in Lemma~\ref{rl:evaluation}; its length is at most
$\dim_{\Omega_j}H^1(C_j\otimes\Omega_j)+\operatorname{rank}M$.
These lengths are bounded.  Include in the integral diagrams the
evaluation map to its two-term complex and a contraction of its cone
to a complex starting in degree one.  Such a contraction uses only
unit pivots over $\mathcal O_j$ and has bounded graded sizes.
All its matrices, including inverses and homotopies, are integral
and pass to the quotient.  Its cone still starts in degree one
after extension to $L_\infty$.  This proves injectivity just as in
Lemma~\ref{rl:evaluation}.

Let $H_j$ be the inner exact kernel.  The product of the central
elements of Lemma~\ref{rl:kernel} is central modulo $\prod_jH_j$.
Its scalar is the same fixed $u^4$ in every component, and
$u^4-1$ remains nonzero in $\mathcal O_\infty$.  The explicit
central-scalar calculation and inflation--restriction give the
claimed restriction injection.  Additional fixed finite diagrams
are treated by the same entrywise operations.
\end{proof}

\begin{lemma}[Irreducibility on the characteristic-zero tests]
\label{rl:irreducibility}
On an inner or outer characteristic-zero test where $\Psi^2\ne1$,
the four-dimensional representation induced from $T\chi\Psi$ is
absolutely irreducible.  If $\Psi^2=1$ on an inner test, it splits
as two inequivalent absolutely irreducible quadratic twists.  In
the outer test $\Psi$ has infinite order.
\end{lemma}

\begin{proof}
The two induced summands are distinguished by $r$-inertia, whose
scalars are $\Psi$ and $\Psi^{-1}$.  On either summand, multiplying
the Tate matrices by scalars does not alter invariant subspaces.
The Tate representation restricted to $G_K$ is absolutely
irreducible under the exclusions already specified.  These facts
give absolute irreducibility of the induction when the two inertia
characters differ.

For the outer argument this assertion requires a uniform finite
witness.  The images of the quadratic subgroups have index at most
two in the fixed compact $2$-adic Tate image.  A compact $2$-adic
analytic group is topologically finitely generated and hence has
only finitely many closed subgroups of index at most two.  Choose
one on an ultrafilter-large set.  Its absolutely irreducible
representation has finitely many matrices spanning
$M_2(\mathbf Q_2)$; the determinant expressing this spanning
property is a fixed nonzero $2$-adic number.  Every one of these
Tate matrices is realized by choices of Galois elements in each
component.  Their accompanying character scalars are units and
do not change the spanning assertion.  Its determinant remains
nonzero in $L_\infty$.  This proves absolute irreducibility in
that field.  Also $1+t$ has infinite order already in every
$\Omega_j$: for a nonzero integer $n$, the first nonzero term of
$(1+t)^n-1$ occurs in finite degree.  It remains a unit modulo
$2$ in every component, so no such power becomes one in the
outer field.

If $\Psi^2=1$, extension across conjugation gives the two quadratic
twists.  They are inequivalent, since a quadratic self-twist of the
Tate representation here would make its restriction to $G_K$
reducible over an algebraic closure.  In the CM case this would
force $K$ to be the CM field, which was excluded.
\end{proof}

\begin{lemma}[Simultaneous finite Chebotarev realization]
\label{rl:chebotarev-evaluation}
For any element $u$ of an exact kernel in Lemma~\ref{rl:kernel},
derivative primes can be chosen, at increasing finite precision,
whose rational Frobenius is $\tau u$ on all retained data.  Their
finite localization is evaluation on $(\tau u)^2$.  The assertion
is compatible with the outer construction of
Proposition~\ref{rl:outer-evaluation} and with any fixed finite
list of previous derivative operations.
\end{lemma}

\begin{proof}
At one Artin stage, a coordinate of a cochain in the image of a
finite model is a locally constant function.  The finitely many
coordinates being used factor through one finite quotient of the
absolute Galois group.  Enlarge it to contain the Tate action at
the required precision, the finitely many ring-class fields and
quadratic characters currently in use, and all old finite data.
Chebotarev realizes the conjugacy class of $\tau u$ on this finite
quotient, with a chosen place realizing the desired representative.
It gives a rational prime inert in $K$, outside the old finite set.
Because $u$ is exactly Tate-trivial, its Frobenius tends to
$\rho_T(\tau)$; in particular $a_\ell\to0$ and $\ell\to-1$.
Because $u$ is exactly ring-class trivial, the square Frobenius on
the $K$-prime is the identity on every retained ring-class field.
The unramified local coordinate is therefore the stated
Frobenius-square evaluation.

Repeat at each finite stage with increasing precision.  All
conditions remain finite, and no density bound or prime-size
bound is required.  For an outer sequence do this separately
for each inner sequence, retaining the relevant finite evaluation
matrices and degree-two identities before taking the outer
quotient.  Approximate inner cycles cause no problem, by the
integrality calculation in Proposition~\ref{rl:outer-evaluation}.
Previously imposed finite diagrams are unchanged.
\end{proof}

In induced-module notation, restriction to the exact kernel is additive.
For a crossed cocycle $b$ and an exact-kernel element $u$,
\[
 b((\tau u)^2)=(1+\tau)b(u).
\]
Indeed $\tau^2=1$, $u$ acts trivially on the coefficients, and the
cocycle identity cancels the two occurrences of $b(\tau)$.  Projection
to either induced summand identifies $\operatorname{im}(1+\tau)$
with a two-dimensional coefficient module, also in characteristic two.

\begin{theorem}[The limit and evaluation interface]
\label{rl:limit-evaluation-package}
For the ring data fixed at the start of this section, retain the finite
target-side arithmetic diagrams with their maps to original cochains, homotopies,
and dualities as in Lemma~\ref{rl:finite-limits}.  The moving characters,
local conditions, and evaluation constructions then have the following
simultaneous properties:
\begin{enumerate}
\item Every specified finite target-side arithmetic diagram admits a bounded free
model together with maps to its original cochains; the specified
maps, homotopies and dualities survive the inner limit.
\item At the stated height-one tests the relaxed active local
complexes are acyclic, and the Selmer duality pairs degrees one
and two.  Integral Kummer conditions at $2N$ include their full
derived degree-zero invariants.
\item The residual sizes needed for the outer construction are
bounded independently of the active conductor support, with
dependence only on $E$, the bound on $\omega(\operatorname{disc}K)$,
and the number of retained derivative primes.
\item Inner and outer $H^1$ are detected by actual crossed
cocycles, including on the exact common kernel.  The evaluations
and their degree-two identities can be realized by finite
Chebotarev conditions without disturbing earlier finite data.
\end{enumerate}
These assertions do not assert an Euler-system reciprocity law or
an integral derivative class that has not first been constructed
at finite level.  Once those finite arithmetic diagrams are
constructed, this Theorem supplies their limit compatibility.
\end{theorem}

\begin{proof}
Combine Lemmas~\ref{rl:finite-model} and~\ref{rl:finite-limits},
Lemma~\ref{rl:kummer}, Proposition~\ref{rl:characters},
Lemmas~\ref{rl:local-acyclic} and~\ref{rl:residual-bound},
Lemmas~\ref{rl:evaluation} and~\ref{rl:kernel},
Proposition~\ref{rl:outer-evaluation}, and
Lemma~\ref{rl:chebotarev-evaluation}.  Every passage uses actual
integral matrices of bounded size; arbitrary evaluations are
retained by their original contraction maps, not by an inference
from a finite list of identities.
\end{proof}
